% from ega4-18-fr.tex
\oldpage[IV-4]{130}

de ses restrictions $g^{(\alpha)}:S'^{(\alpha)}=g^{-1}(S^{(\alpha)})\to S^{(\alpha)}$ le soit, et si l'on pose $S'_0=g^{-1}(S_0)$, $S'^{(\alpha)}_0=(g^{(\alpha)})^{-1}(S^{(\alpha)}_0)$, il y a correspondance biunivoque entre les parties ouvertes et fermées $U$ (resp. $U_0$) de $S'$ (resp. $S'_0$) et les familles $(U^{(\alpha)})$ (resp. $U^{(\alpha)}_0$) où $U^{(\alpha)}$ (resp. $U^{(\alpha)}_0$) est une partie ouverte et fermée de $S'^{(\alpha)}$ (resp. $S'^{(\alpha)}_0$), d'où l'assertion (ii).

\begin{remark}[18.5.7]
\label{IV.18.5.7}
Soient $S=\Spec(A)$ un schéma affine, $S_0$ un sous-schéma fermé de $S$ défini par un idéal $\mathfrak J$ de $A$. Alors, si le couple $(S,S_0)$ est hensélien, l'idéal $\mathfrak J$ est nécessairement contenu dans le \emph{radical} de $A$. En effet, si $\mathfrak m$ est un idéal maximal de $A$, $\mathfrak m$ doit appartenir à $V(\mathfrak J)=S_0$, en vertu de (18.5.4.3), autrement dit on doit avoir $\mathfrak J\subset\mathfrak m$, d'où la conclusion. En particulier, supposons que $S_0$ soit \emph{réduit à un point}, c'est-à-dire que l'idéal $\mathfrak J$ soit \emph{maximal}; alors $\mathfrak J$ doit être le radical de $A$, autrement dit $A$ doit être un anneau \emph{local}, $S_0$ étant l'unique point fermé de $\Spec(A)$.
\end{remark}

\begin{definition}[18.5.8]
\EGAReaderTarget{ega4.c04.def.18-5-8}{IV.18.5.8}
On dit qu'un anneau $A$ est \emph{hensélien} s'il est semi-local et si, en désignant par $\mathfrak r$ le radical de $A$, le couple $(\Spec(A),\Spec(A/\mathfrak r))$ est hensélien. On appelle \emph{schéma local hensélien} un schéma isomorphe au spectre d'un anneau local hensélien.
\end{definition}

\begin{proposition}[18.5.9]
\label{IV.18.5.9}
\begin{enumerate}
  \item[(i)] Pour qu'un anneau semi-local $A$ soit hensélien, il faut et il suffit qu'il soit composé direct d'anneaux locaux henséliens.
  \item[(ii)] Pour qu'un anneau local $A$ soit hensélien, il faut et il suffit que toute $A$-algèbre finie $B$ soit isomorphe à un produit d'anneaux locaux.
\end{enumerate}
\end{proposition}

(i) En effet, la définition, appliquée à $S=\Spec(A)$ et $S_0=\Spec(A/\mathfrak r)$ montre, puisque $S_0$ est un ensemble fini discret et fermé dans $S$, que $S$ est réunion d'un nombre fini de parties ouvertes et fermées $S_i$ ($1\leq i\leq n$) deux à deux disjointes, dont chacune contient exactement un des idéaux maximaux $\mathfrak m_i$ de $A$; la conclusion résulte de (18.5.6, (ii)) et de la remarque (18.5.7).

(ii) Tout morphisme fini $S'\to\Spec(A)$ est de la forme $\Spec(B)\to\Spec(A)$, où $B$ est une $A$-algèbre finie. Si $k$ est le corps résiduel de $A$, $\Spec(B\otimes_A k)$ est un spectre d'anneau artinien, donc fini et discret. Dire que le couple $(\Spec(A),\Spec(k))$ est hensélien signifie donc que $B$ est composé direct d'anneaux $A_i$ tels que $\Spec(A_i\otimes_A k)$ soit réduit à un point, c'est-à-dire que $A_i$ (qui est une $A$-algèbre \emph{finie}) ne doit avoir qu'un seul idéal maximal (Bourbaki, \emph{Alg. comm.}, chap. V, \textsection~2, n\textsuperscript{o}~1, prop. 1).

L'étude des anneaux henséliens est donc essentiellement ramenée à celle des anneaux \emph{locaux} henséliens.

\begin{proposition}[18.5.10]
\label{IV.18.5.10}
Si $A$ est un anneau hensélien, toute $A$-algèbre finie $B$ est un anneau hensélien (donc composé direct d'anneaux locaux henséliens (18.5.9)).
\end{proposition}

En effet, si $\mathfrak r'$ est le radical de $B$, l'image réciproque de $\mathfrak r'$ dans $A$ est le radical $\mathfrak r$ de $A$, et tout idéal de $B$ au-dessus d'un idéal maximal de $A$ est un idéal maximal de $B$, donc l'ensemble $V(\mathfrak r')$ dans $\Spec(B)$ est l'image réciproque de l'ensemble $V(\mathfrak r)$ dans $\Spec(A)$. La proposition est alors une conséquence de (18.5.6, (i)).

\begin{theorem}[18.5.11]
\label{IV.18.5.11}
Soient $A$ un anneau local, $\mathfrak m$ son idéal maximal. Les conditions suivantes sont équivalentes:
\begin{enumerate}
  \item[a)] $A$ est hensélien, autrement dit, toute $A$-algèbre finie $B$ est isomorphe à un produit d'anneaux locaux.

